In coding theory, one main question is 
to determine the best value of $d_{\max}$ for a fixed $n$, $k$ and $q$ such that 
a linear $[n,k,d]_q$ code exists.
There are many bounds, both upper and lower bounds. 
An upper bound tells us that no code with $d \ge d_{\max}$ exists.
A lower bound tells us that a code with $d \ge d_{\max}$ exists. 
The command  

\sourcecode{bounds_for_d_given_n15_k6_q2}


gives upper and lower bounds on the optimal minimum distance $d_{max}$ of a $[15,6]_2$ code.
The values of the Gilbert-Varshamov lower bound 
and the Singleton, Hamming, Plotkin and Griesmer upper bounds are computed. 
The output is:
\orbiteroutput{GRAPHICS/LATEX/bounds_for_d_given_n15_k6_q2}
This shows that $5 \le d_{\max} \le 6.$
The command 

\sourcecode{coding_theory_bounds_q2}

produces a table of bounds for binary codes with $n,k \le 12.$
%The second command produces the latex report below:
%\orbiteroutput{GRAPHICS/LATEX/table_of_bounds_n12_q2}

The command 

\sourcecode{GV_n15_k6_d5}

creates a $[15,6,d]_2$ with minimum distance $g \ge 5$ 
using a greedy algorithm based on the proof of the Gilbert-Varshamov bound.
The code that is produced has the following generator matrix:

\orbiteroutput{GRAPHICS/WRAPPERS/code_GV_n15_k6_d5}

To compute the minimum distance of the code, we store the generator  matrix

\sourcecodevariable{CODE_GV_N15_K6}


and apply the following command to determine the weight enumerator of the code:


\sourcecode{GV_n15_k6_d5_weight_enumerator}

The weight enumerator turns out to be

\orbiteroutput{GRAPHICS/LATEX/GV_n15_k6_d5_weight_enumerator}

From this, we see that the code has minimum distance $6,$ 
which is better than predicted.
%# surprise: d = 6

\bigskip
